\section{The fixed eight-term interpolant}\label{ti:sec:fixed8}
\begin{theorem}\label{ti:main}The rational field defined below has support \(\splitcount{4}{4}\) and at least \(108\) hyperbolic limit cycles. There are two exponential ranges with ten cycles each and a linear range with seven in each open quadrant. Its degree \(D_*\) is the maximum sum of its two collected exponents, attained in the \(j=3\) row of \(P\). If \(\Mclr{}\) denotes its clearing factor as below, \(D_*=1+\Mclr{}(3+\nu_3+\alpha_3)\).\end{theorem}
Section~\ref{sup:primitive-fields} gives the full common factor of the six- through ten-term realizations and the actual degree after cancellation.

\subsection{Two independent regular-energy data}
For $D(w)=\sum_{j=0}^7 c_j w^j$, the area moments for the exponential and
linear horizontal potentials and a quadratic vertical energy are
\begin{align}
 J[D](h)&=\sum_j\frac{2c_j}{2j+1}\int_0^h\frac{(h-s)^{j+1/2}}{1-s}\dd s
                       &&(0<h<1),\notag\\
 T[D](h)&=h^{3/2}\cL[D](h),&
 \cL[D](h)&=\sum_j\frac{4c_j h^j}{(2j+1)(2j+3)}\quad(h>0).
 \label{ti:mom}
\end{align}
Indeed the horizontal potentials on $x>0$ are $1-e^{-x}$ and $x$.
Changing variables from $x$ to its potential value and integrating over the
two vertical branches gives \eqref{ti:mom}. The positive factor $h^{3/2}$
allows interpolation using polynomials alone on the linear range.

Write $r=\sqrt{1-h}$ and $y_j=j+1/2$. Define
$H_{ij}=\binom{y_j}{i}/[y_j(y_j-i)]$.
For $j\ge i$, this is the $i$th derivative at $h=1$ of $J[w^j]$ divided
by $i!$, by differentiation of the integral. The noninteger term in its
expansion has order $r^{2j+1}$, as follows by dividing $s^{j+1/2}$
by $s+r^2$ after substitution in \eqref{ti:mom}.
For $l=2s$, define $B_l$ on the powers $w^s,\ldots,w^{2s}$ by
$H_{0:s} B_l=(0,\ldots,0,1)^{\mathsf T}$ on these columns.
For $l=2s+1$, set $[w^s]B_l=1$ and solve on $w^{s+1},\ldots,w^{2s+1}$
so that the same rows vanish. Only $0\le l\le7$ are used.
The indicated rational square systems are nonsingular by elimination, also
by the endpoint determinant in Lemma~\ref{se:det}.

\begin{lemma}\label{ti:kernel}
Let $\cN$ consist of the real polynomials of degree at most seven satisfying
$[w^0]D=[w^1]D=[w^2]D=0$ and $J[D](1-r^2)=O(r^6)$ as $r\downarrow0$.
Then $\cN=\operatorname{span}(B_6,B_7)$. On this space the map
\[
 D\longmapsto (\cL[D](15/16),\cL[D](49/50))
\]
is invertible. Both evaluation points may also be varied independently in
a neighborhood of this pair.
\end{lemma}
\begin{proof}
With degrees zero through two absent, the first possible noninteger term
has order $r^7$. Thus the remaining three conditions are exactly
$H_{i,3:7}D=0$ for $i=0,1,2$. Their minor on columns $3,4,5$ equals
$183616/1444021425$. The kernel has dimension two and contains the
two indicated polynomials of different degrees.
For $(h_1,h_2)=(15/16,49/50)$ in increasing order, substitution gives
\[
 \det (\cL[B_{j}](h_i))_{i=1,2;\,j=6,7}
 =\frac{32428882679001115293}{2883708677590614016000000000}>0 .
\]
Continuity gives the assertion for nearby nodes.
\end{proof}

Consequently one can fix the first three density coefficients, the three
integer upper-end data of the residual after subtracting those powers,
and two regular linear-range values. These eight data determine a
unique degree-at-most-seven density. This independence lets us assign
the two zeros after constructing the endpoint oscillations.

Put $q=1/16$ and define $D_5=\sum_{l=0}^5
(-1)^{l+\lceil l/2\rceil}q^{(6-l)^2-1}B_l$. Choose $u_6,u_7$ by
\begin{equation}
 A=\frac{D_5+u_6B_6+u_7B_7}{32},\qquad
 \cL[A](15/16)=\cL[A](49/50)=0.                                 \label{ti:choice}
\end{equation}
Lemma~\ref{ti:kernel} makes these exact rational specifications. Numerically
$u_6\simeq-14.27534$ and $u_7\simeq89.34633$. All computations below
use \eqref{ti:choice}. Write $c_j=[w^j]A$.

\begin{lemma}\label{ti:signs}
There are ten disjoint brackets in $0<h<1$ on which $J[A]$ changes sign,
and seven disjoint brackets in $0<h<5/4$ on which $T[A]$ changes sign.
On every bracket the absolute derivative exceeds $2^{-200}$.
At either endpoint the absolute moment exceeds $2^{-710}$.
\end{lemma}
\begin{proof}
The signs on each of the two following ordered mark lists alternate,
starting positive. Both lists begin with
$q^{18},q^{10},q^2,1/20$. Append
\[
\begin{array}{ll}
 J:&1/2,\quad 1-2^{-j}\quad (j=10,30,45,62,75,100),\\
 T:&3/4,\quad 19/20,\quad 1,\quad 5/4 .
\end{array}
\]
In particular the normalized linear polynomial already has seven distinct
positive roots by its signs and degree. To certify the full return we
use one narrow bracket inside each consecutive pair. For every integer
$n$ in Table~\ref{tab:focus-centers}, read in base 16, the bracket is
$n2^{-576}\pm2^{-512}$. All lie above $2^{-90}$, and every exponential
bracket has $1-h>2^{-112}$.

Here are the rational enclosure details. For $h\le1/2$, the ratio
$J[w^j]/T[w^j]$ has a positive series with initial term one and
successive ratios $h(2i+2)/(2i+2j+5)$. After term 1024 the tail is
at most $2^{-1024}$. At the upper brackets $v=1-h\le1/16$, and division gives
\[
 \frac{2j+1}{2}J[w^j](h)=
 \sqrt h\sum_{i=0}^j\frac{(-v)^i h^{j-i}}{j-i+1/2}
       +2(-1)^{j+1}v^{j+1/2}\arccos\sqrt v .
\]
Use $\sqrt v\arccos\sqrt v=(\pi/2)\sqrt v-v\sum_{i\ge0}
\binom{2i}i v^i/[4^i(2i+1)]$. The tail of the series after term 256 is
less than $2^{-1024}$. The alternating series for
$8\arctan(1/5)-2\arctan(1/239)$ enclose $\pi/2$ after 240 terms.
Outward rounding of positive terms on the $2^{-1024}$ grid and integer
square roots give the stated endpoint margin for all 17 brackets.
The linear expression only needs rational evaluation and one square root.
The enclosure calculation checks all endpoints.

For the second derivative use $\max|c_j|<48$. At the upper end the
expression is $\sqrt{1-v}\,P(v)+\sqrt v\arccos\sqrt v\,Q(v)$.
The coefficient norms of $P,Q$ and their first two derivatives are
less than $2^{25}$. For $v\le1/2$ the first two derivatives of
$\sqrt v\arccos\sqrt v$ are bounded in magnitude by $2/\sqrt v$ and
$4/v^{3/2}$. For $h\le1/2$ the positive series divided by $T[w^j]$
has derivatives bounded by 4 and 16. These estimates and \eqref{ti:mom}
bound both second derivatives by $2^{220}$ on their brackets.
The secant slope is larger than $2^{-710}/2^{-512}$ in magnitude.
It differs from any derivative there by at most $2^{221}2^{-512}$.
\end{proof}

\begin{table}[htbp]\centering
\caption{Exact numerators \(n\) for the brackets of Lemma~\ref{ti:signs}.
Read the two lines of each row as one base-16 integer, most significant digit first.
Rows J01--J10 give the exponential moments and T01--T07 the linear moments,
in increasing energy within each group. The center is \(n/2^{576}\) and the radius is \(2^{-512}\).}\label{tab:focus-centers}
\small\ttfamily
\begin{tabular}{@{}l@{\quad}l@{}}\toprule
\textrm{Row} & \textrm{Hexadecimal numerator}\\\midrule
J01 & 4fffffdeeb8a56c7560e9ef5d7b5d7d7fbe547d9662a3d04106188ec80b0c166864f4395f\\*
    & 9e74f6833028db89f0379bbdaed366399efe014da3187b23a597a9a3\\[.3ex]
J02 & 254f6324fd338c91c2eab251c87ce0961731302bd2ef68c34ba07c00728469f3f8690dc05\\*
    & 95e08988a9140fa4a5ce394c14646af6ce44e7b4c5de73f58d308143b7ec4a00\\[.3ex]
J03 & 6dd885a6257f4dbe2779a3df5e242fe6fc18edaad5573f03203013182968b16baaef1fbbb\\*
    & 2798d2a5e3590292f57ae5b19da897c765e5df8462c115c2b6fa82d93ed29bb4805c62\\[.3ex]
J04 & 7e2af0bfb5ae34696f86a033e2676eb98c58f2d5a3d5b9928203cb8ccb43f19c42b43ebfa\\*
    & 1e629dab8b5c87d6de426b9b93f506c4fa2af7a3f5609ce7fb17c1f11cae813da6e2aa7\\[.3ex]
J05 & fdf518a089a18550831cdc41d679dbdd4e5dcbd66c2d16100250951f10aa4c46340b27f7d\\*
    & 08e3d7427ce33f7c09ce794934ecad9f38adb64d047f2c1841f96435a04955ee47e6949\\[.3ex]
J06 & ffffff607e2edbd0e8acc4bc217cedbfd18d79005c6c7f8e33dc2c6087949f87c7a91ac74\\*
    & 02c19c10626dad6690db08a418442fee928d8bba7563fb8cc72a75d4cd4e08384bbbeaa\\[.3ex]
J07 & fffffffffb8765b51aec8b93a01fc3770b8dd813b17901f9ec7469d7ecff5c703d76fe5b6\\*
    & 234445817edf0da5f2c5d36c440555710ce63163618f47b1d092622fe1c6a5d9784bc3e\\[.3ex]
J08 & ffffffffffffffcbb7cc09bc6cfa497b3925cecaada1d2c0b972069309977779401a8eb47\\*
    & 0a18a7a914fa0b7f9e4cef398df4305a2446672631b0139a17349bcc35fce8aab6a96b1\\[.3ex]
J09 & ffffffffffffffffd3dc857a7a9d516893da6da2f098cfafcfdb2200c694f43d7728f1349\\*
    & dd95306ad4d83fde994357bdb91e527842364dc9b9819982908fdcbe9dd828f589331cc\\[.3ex]
J10 & fffffffffffffffffffffff983a008ec5c452f18e381d549125a7fe764126b99ea6d04eb0\\*
    & 25aedebe9d67aeffe95edd90a65a6f708e4b388efa1227a5b8876e5c6e4b6a85d0fbec5\\[.3ex]
T01 & 4fffffdeeb8a56c72857c3adc38731db6b73d227d1b7af728e71ef822c6c3f814de4942a8\\*
    & b61ea7ebce67fb2512c7cfd791c97eae1a2f0a6e770326481b7a5092\\[.3ex]
T02 & 254f6324f7ad68ba9e058ebda3919d50a83fb5aeb87212ccfd659eb32df268152c8457fef\\*
    & 83182e1b43d698260cec4b90c92336e9713174500efdfe20d1fdf590b7d4b101\\[.3ex]
T03 & 6db859566a05e869e3c57a64a4d449c548b62acfc84eb3ce2b15524f6418aa9b3ee73b813\\*
    & 91b5e69ace39a755a44d290b5db5e6399157597caab924bd0d77cb603ca1e2d3dbee36\\[.3ex]
T04 & 79db34e1c370bc4a1d10b043500346aae99c86df5d552cf44441e095f96264b9abb2e34d1\\*
    & a19365c5ee2f4ac36b5dec06125d6f4f91f723ade1c8d25de75651b0789c595799573d5\\[.3ex]
T05 & f000000000000000000000000000000000000000000000000000000000000000000000000\\*
    & 00000000000000000000000000000000000000000000000000000000000000000000000\\[.3ex]
T06 & fae147ae147ae147ae147ae147ae147ae147ae147ae147ae147ae147ae147ae147ae147ae\\*
    & 147ae147ae147ae147ae147ae147ae147ae147ae147ae147ae147ae147ae147ae147ae1\\[.3ex]
T07 & 119d4ecdcc1797a21912f6a89656292c4fa3852a75ea16d7d601b2cc985492c21a946b354\\*
    & dbbe32e2a9d101d97c2758b8a4e9bcd5f5d296903136b2a08364aefecfa949e675fa781f\\[.3ex]
\bottomrule\end{tabular}
\end{table}


\subsection{The eight monomials and their parameters}
Set
\[
 t=\delta=2^{-8192},\quad \tau=t^2,\quad \lambda=t^{16},\quad
 \varepsilon=t^{1024},\quad b_i=\delta^i\ (1\le i\le3),\quad b_H=b_1+b_2+b_3 .
\]
Use the reference potential $K_*=-15/29$ and determinant-ratio slope
$s_*=-76531/180$ from the row calculation below, and put
$\kappa=c_6/(s_*c_7)$, $K=K_*+\kappa\tau$, $a=-(1+3K)$.
The seven controls $X=(a_0,a_2,\sigma,b,\eta_0,\eta_2,\eta_3)^{\mathsf T}$
are the value at $\tau$ of a vector polynomial of degree at most seven,
specified below. Define
\begin{align*}
 C&=(1+2K-\lambda\sigma,-1,-3K+\lambda\sigma,K),&
 \gamma&=(2+\lambda a_0,0,-3+\lambda a_2,1),\\
 \nu&=(\lambda\eta_0,1,2+\lambda\eta_2,3+\lambda\eta_3),&
 g_0&=at^2 b_3/2.
\end{align*}
Set $\alpha_j=\varepsilon\gamma_j/C_j$ and $\beta=\varepsilon\lambda b/g_0$.
Let $\Mclr{}$ be twice the least common multiple of the reduced denominators of
all $\alpha_j,\nu_j,b_i,\beta$. The collected polynomial is
\begin{align}
 P(x,y)&=\sum_{j=0}^3 C_j x^{1+\Mclr{}(1+\alpha_j)}y^{\Mclr{}(2+\nu_j)},\notag\\
 Q(x,y)&=\tfrac12 at^2\bigl(-b_Hx^{2\Mclr{}}y^{1+\Mclr{}}
 +b_1x^{\Mclr{}(1+b_1)}y^{1+\Mclr{}}\notag\\
 &\hspace{24mm}{}+b_2x^{\Mclr{}(1+b_2)}y^{1+\Mclr{}}
 +b_3x^{\Mclr}y^{1+\Mclr{}(1+\beta)}\bigr).                               \label{ti:poly}
\end{align}
In logarithmic coordinates $u=\Mclr{}\log|x|, v=\Mclr{}\log|y|$, with $L=e^v$, the
field before the positive time change is
\[
 \dot u=\sum_j C_j e^{\alpha_j u} L^{1+\nu_j},\qquad
 \dot v=\tfrac{at^2}{2}(-b_H e^u+b_1 e^{b_1u}+b_2 e^{b_2u})+g_0e^{\beta v}.
\]
Multiplication of the pullback by $\Mclr{}x^{\Mclr}y^{\Mclr}$ gives \eqref{ti:poly} in every
open quadrant. The coefficient bounds below imply $1/2<a<1$,
$1/40<|C_j|<4$, $|\nu_j-j|<1/4$ and $|\alpha_j|,|\beta|<1/4$.
All exponents are nonnegative integers. The four $y$ powers in $P$ are
different. The four $x$ powers in $Q$ are different since $1>b_1>b_2>0$.
The coefficients have thus survived separate collection.

For completeness we give a rational row formula for $X$. Put $z=L-1$,
$F_0=az+Kz^3$, $g=(z^2-3)/(1+3K-Kz^2)$,
$d=(2,0,-3,1)$ and $C^0=(1+2K,-1,-3K,K)$. Write $H_j(z)=\ev(L^j\log L)$,
where $\ev$ takes the even part in $z$, and
$\mathcal U f=(g(z)\int_0^z f(s)\dd s)'$. Define $M_i(K)$ as the row
of $z^{2i}$ coefficients of
\[
 1,\quad\ev L^2,\quad\mathcal U(\ev(L^2-1)),\quad\ev(1/L),\quad
 d_jH_j+C^0_j\mathcal U H_j\quad(j=0,2,3).
\]
For $0\le i\le7$, set
\[
 G_i(K)=\sum_{j=0}^i\binom{-(2i+1)/2}{i-j}(K/(2a))^{i-j}M_j(K),
 \qquad P_i(\tau)=(1+3K)^{i+1}G_i(K).
\]
The $P_i$ are row polynomials of degrees at most $i+1$. Determine $X$
by coefficient comparison, with $S=(0,1,2,3,4,5,7)$, from
\begin{equation}
 P_i(\tau)X(\tau)=(1+3K)^{i+1}c_i\tau^{7-i}\pmod{\tau^8}
 \qquad (i\in S).                                                    \label{ti:shape}
\end{equation}
The constant selected matrix is invertible. For the $M_i$, its unscaled
order-seven minor on $S$ at $K_*$ is $20511149/833191988428800$.
The minor on the first seven rows has a simple zero there.
At $K_*$ the lower triangular transformation to $G_i$ preserves the
derivative $s_*$ of the ratio of the minor on rows $0,\ldots,6$ to the minor on $S$.
With the chosen $\kappa$, the omitted row $i=6$ of \eqref{ti:shape} holds modulo $\tau^2$.
The residual calculation on p.~\pageref{ti:row-residuals} verifies this congruence directly.

\subsection{From the brackets to hyperbolic cycles}
At $\varepsilon=0$ the logarithmic system is Hamiltonian with
$H=at^2W(u)/2+V(z)$, where
\[
 W(u)=b_H(e^u-1)+(1-e^{b_1u})+(1-e^{b_2u})-b_3u,\qquad
 V(z)=\int_0^z\sum_j C_j(1+s)^{\nu_j}\dd s .
\]
On the annuli in use $Y=z\sqrt{2V/(az^2)}$ is increasing with $Y(0)=0$.
The scaled energy is $E=2H/(at^2)=W+Z^2$, $Z=Y/t$. In
$x_i=-b_i u>0$ the limiting horizontal potential is $i-1+1-e^{-x_i}$
for $i=1,2$ and $2+x_i$ for $i=3$. Thus $h=E-(i-1)$ in \eqref{ti:mom}.
The negative $x_i$ side tends to a confining wall.

Write $M(E)$ for the first variation in $\varepsilon$ of the return
displacement in $H$. Green's theorem on the clockwise oval gives the
sublevel integral with density
$D_\lambda(z)=\sum_j\gamma_jL^{\nu_j}+\lambda b/L$ in $\dd u\dd z$.
In the variable $Y$ only its transported even part
$S_\lambda(Y)=\ev_Y[D_\lambda(z(Y))z'(Y)]$ contributes.
At $\lambda=0$ the density is $z^3-3z$ and $Y$ is odd.
Differentiating at fixed $Y$ moves $z$ by $-\partial_\lambda V/F_0$.
Differentiating the density and its Jacobian gives the rows $M_i$
above, followed by the change from $z$ to $Y$ giving $G_i$.
The matching in \eqref{ti:shape} therefore applies to the first
variation of this even density.

We apply the finite estimates of Section~\ref{fin:cert}
at the same values of $t,\delta,\varepsilon$. We detail their hypotheses
for the present density because the bracket centers and controls have
changed. For $R_i=P_iX-(1+3K)^{i+1}c_i\tau^{7-i}$,
exact rational multiplication and elimination give
\[
 |\kappa|<1,\qquad
 \max_j\sum_l|[\tau^l]X_j|<535481988140936,\qquad
 \max_i\sum_l|[\tau^l]R_i|<14030792279473 .
\]
Both integer bounds are below $2^{60}$. \phantomsection\label{ti:row-residuals}The first eight residual
coefficients vanish on $S$, and the first two vanish in row 6.
The rational solve checks these statements directly.
On $0\le w\le32$ the coefficient sums bound both $|A(w)|,|A'(w)|$
by $2^{48}$. All brackets satisfy the regular-energy restrictions of the
finite estimates by Lemma~\ref{ti:signs}. These are the density-dependent
hypotheses of the transfer.

More explicitly, replacing the finite horizontal potential by its model
costs less than $2^{-4000}$ in value and $2^{-770}$ in derivative of the
normalized area integral. This follows by splitting off
$|x_i|\le\sqrt\delta$ and using $|W-W_*|<512\delta$ on the remaining
sublevel, with a further split at distance $2\delta^{1/4}$ from the
turning energy for the derivative. The vertical estimate
$|S_\lambda(tZ)/(\lambda t^{14})-A(Z^2)|<2^{230}t^2$ for $|Z|\le6$
uses the row residual orders and the analytic inverse in $|Y|\le1/16$.
Integration for $b_iM(E)/(\lambda t^{15})$ costs less than $2^{-15000}$
in value and derivative.

Finally let $\Pi$ be the actual one-turn return in $E$ on the positive
horizontal ray of $(\operatorname{sgn}(u)\sqrt W,Z)$.
On $2^{-95}\le E\le32$ the angular denominator stays positive and the
energy equation is $dE/d\theta=\varepsilon\mathcal F(E,\theta,\varepsilon)$.
The function and its derivatives of orders $E,EE,\varepsilon,E\varepsilon$
are bounded by $N_0^{200}$, $N_0=1/t$, as in Section~\ref{fin:cert}.
Those bounds use $|\gamma_j/C_j|<256$ and $|\lambda b/g_0|<1$,
which follow from the displayed control bounds. Integrating this equation
and its first variational equation for one turn gives
\[
 \left\|\frac{b_i a}{2t^{29}}\frac{\Pi(E)-E}{\varepsilon}
       -\frac{b_i M(E)}{\lambda t^{15}}\right\|_{C^1}
       \le256\varepsilon N_0^{429}<2^{-1000},
\]
where the norm is the maximum of the value and derivative suprema on
each bracket. Combined with the preceding estimates this preserves
the two endpoint signs and the strict derivative sign. Each bracket gives
a simple return fixed point. Its multiplier is positive by the scalar
variational equation and differs from one. The intervals are disjoint on
one ray and the angle is strictly increasing, so distinct fixed points
give distinct periodic orbits. The quadrant diffeomorphisms and their
positive time factors preserve the return multipliers. This proves
Theorem~\ref{ti:main}.

The two equalities in \eqref{ti:choice} specify exact model energies. For
the full polynomial the corresponding energies lie in the certified
brackets translated by two. All strict estimates persist for sufficiently
close pairs of interpolation nodes. Rational pairs in that neighborhood
give a family of rational polynomial fields by the same construction.
The pair in Lemma~\ref{ti:kernel} has the stated finite certificate.
